\documentclass[12pt]{article}
\usepackage[T1]{fontenc}
\usepackage[utf8]{inputenc}

\usepackage{geometry}
\geometry{a4paper, margin=2.5cm}
\usepackage{setspace}
\setlength{\parindent}{0pt}
\setlength{\parskip}{0.8em}

\begin{document}

\thispagestyle{empty}

\noindent
The Editors\\
\textit{Philosophy \& Technology}\\
Springer Nature

\bigskip
\noindent
Independent Researcher\\
Kanagawa, Japan\\
lucaremm@gmail.com\\
\today

\bigskip

\noindent Dear Editors,

\bigskip

I am pleased to submit for your consideration the manuscript titled \textbf{``Reconfiguration of Inquiry in LLM Environments: Origin Opacity and the Conditions of Human Responsibility''} for publication as a regular article in \textit{Philosophy \& Technology}.

\bigskip

\noindent\textbf{Central contribution}

This paper introduces and theorizes \textit{origin opacity}: a structural condition in which the pre-reflective generative structure of human inquiry is progressively reconfigured through repeated interaction with large language model (LLM) environments, in a manner that is not sufficiently visible to the subject.

The paper makes three interconnected contributions. First, it identifies a qualitatively new layer of technological transformation that extends beyond the classical frameworks of Wiener, Ellul, and Arendt: whereas prior technology philosophy addressed the transformation of human activity, judgment, and responsibility conditions, LLM environments can penetrate the pre-reflective generative structure of inquiry itself. Second, drawing on Polanyi's structure of tacit knowledge, Merleau-Ponty's phenomenology of the silent logos and chiasm, and Gibson's ecological theory of affordances, the paper provides a theoretically grounded account of why this compression remains structurally imperceptible to the subject. Third, it argues that this opacity deepens rather than dissolves the asymmetry of responsibility: the origin conditions of inquiry are impaired while terminal responsibility remains with the human being. The paper therefore reformulates irreversibility, fixation of responsibility, and terminal human responsibility not as supplementary ethical requirements but as prior conditions for technological progress to remain humanly bearable.

\bigskip

\noindent\textbf{Relation to existing literature}

While adjacent discussions address the transformation of the information environment (Floridi), the extension of cognition through AI tools (Heersmink), and AI ethics more broadly (Coeckelbergh, Gabriel), no existing work, to my knowledge, has focused specifically on the pre-reflective generative structure of inquiry as a site of LLM-induced transformation. The paper situates itself within the classical trajectory of Wiener, Ellul, and Arendt, while extending that trajectory to a qualitatively new layer that their frameworks did not address. The paper further engages the meaningful human control framework of Santoni de Sio and van den Hoven, arguing not against it but for its prior condition: that the generative structure of inquiry from which such control must arise may itself be impaired.

\bigskip

\noindent\textbf{Appropriateness for \textit{Philosophy \& Technology}}

The paper is addressed directly to \textit{Philosophy \& Technology}'s mandate of examining ``the conceptual nature and practical consequences'' of technology. It engages the philosophy of technology tradition (Ellul, Arendt), phenomenology (Merleau-Ponty, Polanyi, Gibson), and responsibility theory (Santoni de Sio and van den Hoven, EU AI Act), while making an original theoretical contribution that bears on urgent practical questions of AI governance and institutional design. The paper is submitted as an independent researcher, a position that, as the paper itself argues, may enable engagement with questions that institutionally situated researchers may find structurally difficult to pursue.

\bigskip

\noindent\textbf{Methodological transparency}

In keeping with the paper's own argument, it should be disclosed that the author used LLM-assisted tools during the preparation of this manuscript for structural refinement, linguistic elaboration, and editorial organization. Conceptual development, theoretical claims, and all final decisions regarding content were made by the author. Dialogue logs and observation records have been preserved and can be disclosed upon request. These records are preserved not for the purpose of empirical verification but to ensure methodological transparency.

\bigskip

\noindent\textbf{Submission statement}

This manuscript has not been published previously and is not under consideration for publication elsewhere. It presents original research for which the author bears full responsibility.

\bigskip

I would be grateful for the opportunity to have this work considered by \textit{Philosophy \& Technology}. I look forward to your response.

\bigskip

\noindent Sincerely,

\bigskip\bigskip

\noindent Madoka Masui\\
Independent Researcher\\
Kanagawa, Japan\\
lucaremm@gmail.com

\end{document}
